% ============================================================
%  H. Høffding, Psychology in Outline on the Basis of Experience
%  2nd, partly revised edition.
%  Copenhagen: P. G. Philipsen, 1885. IV + 421 pp.
%  TRANSLATION (English)
% ============================================================
\documentclass[12pt,a4paper]{book}
\usepackage[utf8]{inputenc}
\usepackage[T1]{fontenc}
\usepackage[english]{babel}
\usepackage{microtype}
\usepackage{setspace}
\usepackage[margin=1.2in]{geometry}
\usepackage{fancyhdr}
\setlength{\headheight}{28pt}
\addtolength{\topmargin}{-13.5pt}
\usepackage{hyperref}

\hypersetup{
  pdftitle={Psychology in Outline},
  pdfauthor={Harald Høffding},
  hidelinks
}

\pagestyle{fancy}
\fancyhf{}
\fancyhead[LE,RO]{\thepage}
\fancyhead[RE]{\textit{Psychology in Outline}}
\fancyhead[LO]{\textit{\leftmark}}

\setstretch{1.3}

\title{\textbf{Psychology in Outline}\\[6pt]
  \large on the Basis of Experience\\[6pt]
  \normalsize\textit{Psykologi i Omrids paa Grundlag af Erfaring}\\[4pt]
  Second, partly revised edition}
\author{Harald Høffding\\[4pt]
  \small Translated from the Danish}
\date{Copenhagen: P.\ G.\ Philipsen, 1885}

\begin{document}
\maketitle
\thispagestyle{empty}
\newpage

\tableofcontents
\newpage


\chapter{The Subject Matter and Method of Psychology}
\label{ch:subject}
% [text to be added]

\chapter{Soul and Body}
\label{ch:soul-body}
% [text to be added]

\chapter{The Conscious and the Unconscious}
\label{ch:conscious-unconscious}
% [text to be added]

\chapter{Classification of the Psychological Elements}
\label{ch:classification}
% [text to be added]

\chapter{The Psychology of Cognition}
\label{ch:cognition}

\section{Sensation}
\label{sec:sensation}
% [text to be added]

\section{Representation}
\label{sec:representation}
% [text to be added]

\section{The Apprehension of Time and Space}
\label{sec:time-space}
% [text to be added]

\section{The Apprehension of the Real}
\label{sec:real}
% [text to be added]

\chapter{The Psychology of Feeling}
\label{ch:feeling}

\section{Feeling and Sense-Sensation}
\label{sec:feeling-sensation}
% [text to be added]

\section{Feeling and Representation}
\label{sec:feeling-representation}
% [text to be added]

\section{Egoistic and Sympathetic Feeling}
\label{sec:egoistic-sympathetic}
% [text to be added]

\section{The Physiology and Biology of Feeling}
\label{sec:feeling-physiology}
% [text to be added]

\section{The Validity of the Law of Relation for Feelings}
\label{sec:relation-law-feeling}
% [text to be added]

\section{The Influence of Feeling on Cognition}
\label{sec:feeling-cognition}
% [text to be added]

\chapter{The Psychology of Will}
\label{ch:will}

\section{The Originality of Will}
\label{sec:will-originality}
% [text to be added]

\section{Will and the Other Elements of Consciousness}
\label{sec:will-consciousness}

\subsection{The Higher Development of Will, Conditioned by the
  Development of Cognition and Feeling}
\label{ssec:will-higher}
% a. The Psychology of Drive.
% b. The Wish.

% [text to be added: a. The Psychology of Drive, b. The Wish]

% c. Deliberation, Resolve, and Decision.

When the momentary impulses rule, there is properly no inner center, no
self in consciousness. But the more memory and the ruling feelings
determined by it assert themselves, the more does the human being's
whole nature, not merely a single, momentarily predominant side of it,
gain influence over action. The human being's true self has its
expression in the thoughts and feelings that in the course of life have
taken deepest root in him. (Cf.~V B, 5). And only when action has come
to be determined by this firm core can the human being be said to have
willed his action, to have determined himself. Will is to that extent
determined by the ruling passion as the expression of the result that
the human being's development has deposited up to the moment of action.
Now it is the elements of cognition, now the elements of feeling, that
predominate, so that one can distinguish between a thought-will and a
feeling-will. We must here remember that it is not always the most
vehement feeling of the moment that is in reality the strongest---that
is, that has the deepest root in the individual's nature. All the more
important is it that the moment should not be the sole determining
factor; therefore the formation of the aforementioned interval is so
important. It is filled by deliberation, the inner debate in
consciousness, in which all representations and feelings psychologically
possible for the individual can speak a word. Through this debate the
mere resolve becomes decision, and a choice is made among the
possibilities that present themselves. The dice are cast.

This is the highest form of willing. It is predominantly determined from
within, not by the single sensation. Drive knows only one possibility,
one motive; genuine will unfolds itself through the interaction or
struggle of several motives and possibilities. It is often determined by
something that lies far beyond the moment, indeed beyond the
individual's entire reality, but which nonetheless asserts itself in the
individual's consciousness, is represented within it. The psychological
understanding of the willed action depends on whether it is possible to
trace the individual's course of development and follow it step by step
to the moment of action, each individual transition being explained
according to the general psychological laws. This is of course an
ideal; but it indicates the principle and the presupposition of all
psychological inquiry and of all practical dealings among human beings.

\subsection{The Counter-Effect of Will on Cognition and Feeling}
\label{ssec:will-counter}
% [text to be added]

\subsection{The Relation of Opposition between Will and the Other
  Elements of Consciousness (Concentration and Expansion)}
\label{ssec:concentration-expansion}
% [text to be added]

\subsection{The Consciousness of Will}
\label{ssec:consciousness-will}
% a. The consciousness of willing.
% b. The problem of reality in the domain of will.
% [text to be added]

\subsection{Will and the Unconscious Life of the Soul}
\label{ssec:will-unconscious}
% a. The center of consciousness not always the center of individuality.
% b. Determinism and indeterminism.
% [text to be added]

\section{The Individual Character}
\label{sec:character}
% [text to be added]

\end{document}
